\appendices
\section{Additional Tables}
These tables hold detail that the main text summarises with the numbers in \texttt{numbers.tex}. The B5 tables are post-hoc exploratory: B5 and D1 were designed after the test results were known.

\begin{table*}[!tb]
\caption{Share of qualified students inside the top-20 pool, by resume style (counts in brackets), pooled over the test requirements.}
\label{tab:archetype}
\centering\small
\resizebox{\textwidth}{!}{\begin{tabular}{lcccccccc}
\toprule
Method & standard & paraphrased & implied & abbreviated\_noisy & sparse\_strong & ineligible\_strong & generalist\_average & stuffed\_weak \\
\midrule
B1 ATS strict & 78\% (60/77) & 0\% (0/51) & 44\% (8/18) & 14\% (3/21) & 64\% (7/11) & 100\% (1/1) & 62\% (5/8) & 0\% (0/3) \\
B1b ATS lenient (>=50\%) & 92\% (71/77) & 0\% (0/51) & 56\% (10/18) & 38\% (8/21) & 100\% (11/11) & 100\% (1/1) & 38\% (3/8) & 67\% (2/3) \\
B1c ATS + synonyms & 94\% (72/77) & 0\% (0/51) & 56\% (10/18) & 90\% (19/21) & 100\% (11/11) & 100\% (1/1) & 62\% (5/8) & 0\% (0/3) \\
B2 BM25 & 99\% (76/77) & 2\% (1/51) & 94\% (17/18) & 43\% (9/21) & 82\% (9/11) & 100\% (1/1) & 75\% (6/8) & 0\% (0/3) \\
B3 Dense only & 61\% (47/77) & 6\% (3/51) & 78\% (14/18) & 52\% (11/21) & 82\% (9/11) & 100\% (1/1) & 50\% (4/8) & 0\% (0/3) \\
B3b Hybrid (RRF) & 94\% (72/77) & 4\% (2/51) & 94\% (17/18) & 57\% (12/21) & 100\% (11/11) & 100\% (1/1) & 88\% (7/8) & 0\% (0/3) \\
B4 Proposed (K=20) & 61\% (47/77) & 6\% (3/51) & 78\% (14/18) & 52\% (11/21) & 82\% (9/11) & 100\% (1/1) & 50\% (4/8) & 0\% (0/3) \\
\bottomrule
\end{tabular}
}
\end{table*}

\begin{table*}[!tb]
\caption{Post-hoc exploratory: B5 (hybrid retrieval plus the same re-ranker), mean with 95\% bootstrap CI over requirements (dev and test shown separately).}
\label{tab:b5}
\centering\small
\resizebox{\textwidth}{!}{% POST-HOC EXPLORATORY (designed after seeing the test results)
\begin{tabular}{lcccccc}
\toprule
Method (post-hoc) & nDCG@10 & MRR & P@10 & Recall@10 (capped) & FRR@10 & FRR@20 \\
\midrule
B5 hybrid + LLM re-rank, dev (n=5) & 0.724 [0.574, 0.838] & 1.000 [1.000, 1.000] & 0.740 [0.540, 0.920] & 0.740 [0.540, 0.920] & 0.260 [0.080, 0.460] & 0.393 [0.290, 0.477] \\
B5 hybrid + LLM re-rank, test (n=10) & 0.772 [0.707, 0.824] & 1.000 [1.000, 1.000] & 0.860 [0.790, 0.930] & 0.860 [0.790, 0.930] & 0.140 [0.070, 0.210] & 0.314 [0.233, 0.384] \\
\bottomrule
\end{tabular}
}
\end{table*}

\begin{table*}[!tb]
\caption{Post-hoc exploratory: share of qualified students inside the candidate pool, for the dense pool (B4), the hybrid ranking (B3b) and the hybrid pool (B5).}
\label{tab:b5arch}
\centering\small
\resizebox{\textwidth}{!}{% POST-HOC EXPLORATORY (designed after seeing the test results)
\begin{tabular}{lcccccccc}
\toprule
Pool (post-hoc) & standard & paraphrased & implied & abbreviated\_noisy & sparse\_strong & ineligible\_strong & generalist\_average & stuffed\_weak \\
\midrule
B4 (dense pool) & 61\% (47/77) & 6\% (3/51) & 78\% (14/18) & 52\% (11/21) & 82\% (9/11) & 100\% (1/1) & 50\% (4/8) & 0\% (0/3) \\
B3b hybrid ranking (no LLM) & 94\% (72/77) & 4\% (2/51) & 94\% (17/18) & 57\% (12/21) & 100\% (11/11) & 100\% (1/1) & 88\% (7/8) & 0\% (0/3) \\
B5 (hybrid pool) & 94\% (72/77) & 4\% (2/51) & 94\% (17/18) & 57\% (12/21) & 100\% (11/11) & 100\% (1/1) & 88\% (7/8) & 0\% (0/3) \\
\bottomrule
\end{tabular}
}
\end{table*}

\begin{table*}[!tb]
\caption{Name-swap test: control (same-group name) versus cross-group names, with the fairness instruction on and off.}
\label{tab:fair}
\centering\small
\resizebox{\textwidth}{!}{\begin{tabular}{llccc}
\toprule
Condition & Metric & Control & Cross-group & Cross$-$control \\
\midrule
instruction ON & mean abs score change (points) & 3.750 [1.750, 6.350] & 2.383 [1.450, 3.433] & -1.367 [-3.517, 0.250] (not distinguishable) \\
instruction ON & share of swaps with abs change >= 5 & 0.170 [0.100, 0.250] & 0.160 [0.103, 0.223] & -0.010 [-0.073, 0.053] (not distinguishable) \\
instruction ON & mean abs rank shift in 20-pool & 0.710 [0.330, 1.140] & 0.403 [0.220, 0.617] & -0.307 [-0.697, 0.047] (not distinguishable) \\
instruction ON & top-10 boundary-crossing rate & 0.070 [0.020, 0.120] & 0.033 [0.013, 0.060] & -0.037 [-0.083, 0.007] (not distinguishable) \\
instruction ON & mean abs cosine-similarity shift (retrieval) & 0.014 [0.012, 0.016] & 0.014 [0.013, 0.015] & -0.000 [-0.002, 0.002] (not distinguishable) \\
instruction OFF & mean abs score change (points) & 1.850 [0.950, 2.900] & 2.083 [1.233, 3.083] & 0.233 [-0.617, 1.100] (not distinguishable) \\
instruction OFF & share of swaps with abs change >= 5 & 0.150 [0.080, 0.220] & 0.143 [0.090, 0.203] & -0.007 [-0.060, 0.043] (not distinguishable) \\
instruction OFF & mean abs rank shift in 20-pool & 0.580 [0.250, 0.980] & 0.577 [0.283, 0.940] & -0.003 [-0.263, 0.240] (not distinguishable) \\
instruction OFF & top-10 boundary-crossing rate & 0.060 [0.020, 0.110] & 0.063 [0.027, 0.107] & 0.003 [-0.023, 0.030] (not distinguishable) \\
instruction OFF & mean abs cosine-similarity shift (retrieval) & 0.014 [0.012, 0.016] & 0.014 [0.013, 0.015] & -0.000 [-0.002, 0.002] (not distinguishable) \\
\bottomrule
\end{tabular}
}
\end{table*}
\begin{table}[!t]
\caption{Post-hoc exploratory: D1, scoring paraphrased resumes with retrieval bypassed (per requirement).}
\label{tab:d1}
\centering\scriptsize
\resizebox{\columnwidth}{!}{% POST-HOC EXPLORATORY (designed after seeing the test results)
\begin{tabular}{lcccccc}
\toprule
Req & Para. qual. & Non-qual. & Score qual. & Score non-qual. & AUC & Top-10 \\
\midrule
R02 & 4 & 15 & 83.8 & 42.4 & 0.808 & 4/4 \\
R03 & 5 & 15 & 62.0 & 17.1 & 0.947 & 5/5 \\
R05 & 6 & 15 & 95.0 & 32.9 & 0.867 & 6/6 \\
R06 & 3 & 15 & 85.0 & 47.7 & 0.844 & 3/3 \\
R08 & 7 & 15 & 94.3 & 16.7 & 0.995 & 7/7 \\
R10 & 6 & 15 & 83.3 & 35.1 & 0.900 & 6/6 \\
R11 & 3 & 15 & 86.7 & 22.3 & 0.956 & 3/3 \\
R13 & 7 & 15 & 49.3 & 25.7 & 0.767 & 6/7 \\
R14 & 5 & 15 & 100.0 & 32.7 & 0.933 & 5/5 \\
R15 & 5 & 15 & 81.8 & 22.2 & 0.847 & 4/5 \\
\bottomrule
\end{tabular}
}
\end{table}

\FloatBarrier
